\documentclass{article}
\PassOptionsToPackage{numbers}{natbib}
\usepackage[sglblindworkshop]{neurips_2026}
\usepackage[utf8]{inputenc}
\usepackage[T1]{fontenc}
\usepackage{hyperref}
\usepackage{url}
\usepackage{booktabs}
\usepackage{tabularx}
\usepackage{amsfonts}
\usepackage{amsmath}
\usepackage{microtype}
\usepackage{xcolor}
\usepackage{graphicx}
\workshoptitle{Galileo Hackathon 2026}
% Venue notice only: preserve the official typography, margins and spacing.
% This is an independent hackathon report, not a NeurIPS/ML4PS submission.
\makeatletter
\renewcommand{\@noticestring}{Galileo Hackathon 2026. Participant research report.}
\makeatother

% ---------------------------------------------------------------------------
% FILL-IN MARKERS (delete these macros' uses before submitting).
%   \fillin{WHO}{what to write}  -> red, owner-tagged placeholder
%   \tbd                          -> a number we still have to measure/copy
% Find all remaining: grep -n 'fillin\|tbd' report/main.tex
% Owners: PRISHA, ZACH, BOTH. Core judgments (hypothesis, insights, why it
% works, limitations, loophole reasoning) must be written by us, not AI.
% ---------------------------------------------------------------------------
\newcommand{\fillin}[2]{\textcolor{red}{[\textbf{#1}: #2]}}
\newcommand{\tbd}{\textcolor{red}{\textbf{TBD}}}
% Figure that shows a placeholder box until the file exists in report/figs/.
\newcommand{\figorbox}[2]{%
  \IfFileExists{#1}{\includegraphics[width=\linewidth]{#1}}%
  {\fbox{\parbox[c][3.2cm][c]{0.95\linewidth}{\centering\textcolor{red}{Missing \texttt{\detokenize{#1}}\\#2}}}}}

\title{\fillin{BOTH}{Title stating the specific finding, written last}}
\author{Hype Cycle \\ Prisha Maiti\fillin{PRISHA}{affiliation} \quad
Zachary Doroshenko\fillin{ZACH}{affiliation}
\\ Team ID: G2-15 \quad Task slug: \texttt{exoplanet-orbits}}

\begin{document}
\maketitle

\begin{abstract}
\fillin{BOTH}{Write last, 4--6 sentences: the scientific question (recover
period $P$ and semi-amplitude $K$ of a single planet from sparse radial
velocities); the initial idea; the most informative experiment; the final
finding with its main limitation. Quote the hidden-test score and final
submission ID.}
\end{abstract}

% ===========================================================================
\section{Scientific goal \& initial hypothesis}
% Owner: PRISHA (drafts), ZACH (checks physics)
An orbiting planet makes its star wobble around their common centre of
mass; the radial-velocity (RV) method measures the star's line-of-sight
velocity from the Doppler shift of its spectral lines, and remains the main
way to measure planet masses. For one planet, the velocity follows a
Keplerian curve set by the period $P$, the semi-amplitude $K$ (half the
peak-to-peak swing) and the eccentricity $e$ (0 gives a sine wave; high $e$
gives a brief spike at periastron), plus nuisance terms for orientation,
phase and offset. The challenge~\cite{challenge} is the inverse problem:
recover $P$, $K$ and $e$ from a few dozen noisy night-time measurements.
It is hard because six parameters must be fitted from few points; the
noise (reported errors plus unreported stellar jitter) is comparable to the
signal; nightly and seasonal sampling create daily and yearly
\emph{aliases}, false periods that fit as well as the true one; and eccentric
spikes are easily missed by sparse sampling.

\textbf{Task facts (from the official task card, \texttt{TASK\_CARD.md} in the task package~\cite{galileo2026}).} Input: up to 60 observations per star
(time, RV, reported error, mask). Output: $[\log_{10}P, \log_{10}K, e]$.
Metric: \texttt{success\_rate} (higher is better), the share of stars with
$|P_\text{pred}/P_\text{true}-1|<1\%$ and $|K_\text{pred}/K_\text{true}-1|<10\%$;
$e$ is not scored. Groups: in-distribution, high eccentricity
($e$ 0.6--0.9, beyond the training range), and sparse (15 observations).
The official baseline scores 0.4715 on validation; the per-star Fisher
(Cram\'er--Rao) information ceiling is 0.536, and a perfect period search with
the baseline likelihood would reach only 0.479.

\textbf{Why a high success rate is out of reach.}
% Facts with sources (TASK_CARD "Data" and "Difficulty", official analysis/headroom.py,
% results/headroom_*.json). PRISHA/ZACH: rephrase in our own words and keep the caveat.
The limit comes from the data, not from the method. Per-point noise is
comparable to the signal: reported errors are 1--5\,m/s plus an unreported
stellar jitter of 1--3\,m/s, while $K$ is log-uniform in 2--100\,m/s, so
23.7\% of validation stars have $K<5$\,m/s, sampled 15--60 times. The success
box is strict ($P$ within 1\%, $K$ within 10\%). The official Fisher
(Cram\'er--Rao) analysis, which propagates the true noise to the
uncertainty of $(\ln P, \ln K)$ at the true orbit, gives each star a
probability of landing in that box: 20.3\% of validation stars have
probability $<0.1$ (no method can be expected to recover them), and only
31.4\% have probability $>0.9$, where the baseline already succeeds 97.7\%
of the time (our final model: 97.9\%). The mean is 0.536 overall (in-dist.\ 0.657, high-$e$ 0.357,
sparse 0.472). Most alias failures are not search errors either: in only 7
of the baseline's alias failures does the true orbit fit the data better
than the baseline's answer, so the data genuinely prefer the alias at this
noise level. A perfect period search adds only +0.0075 (task card) and
+0.011 for our final model, and four search changes we tried moved 0 of
1,000 stars (R07). Two caveats: the Fisher value is a \emph{benchmark}
rather than a strict maximum. It is local (it assumes the right peak was
found and ignores aliases, which makes it optimistic) and applies to
unbiased estimators (a biased estimator, such as our $K$ correction, can
beat it on individual stars). Also, with 4,000 test stars one binomial
standard error is $\approx$0.8 points, so differences between methods are
small compared with the noise in a single score.
In short, the noise in the data, not the method, caps the achievable
score: an unbiased estimator that always finds the right orbit is expected
to recover only about 54\% of the planets, so the headroom over the baseline
is about 6 points. We therefore
judge our changes by how much of that gap they close, and we measure them
with paired comparisons on the same stars, which separate a real gain of a
few stars from the sampling noise of a single score.

\textbf{Initial hypothesis.} We expected period aliasing to be the main
bottleneck: because the stars are observed only at night and only in
certain seasons, the true period has daily and yearly aliases, and we
expected most failures to come from the method selecting one of them.
A more thorough period search and better alias handling should then raise
the success rate. The hypothesis would be contradicted if a perfect period
search added little, or if the data themselves favoured the alias in most
alias failures.

% ===========================================================================
\section{Research trajectory}
% Owner: ZACH writes rows as runs finish; PRISHA keeps report/EXPERIMENT_LOG.md
% in sync and writes the connecting prose.
The evidence contradicted our initial hypothesis. The task card's analysis
showed that a perfect period search would add only +0.0075, and that in
almost all alias failures the data genuinely prefer the alias; our own
search changes later moved 0 of 1,000 stars (R07). The limiting factor is
the noise in the data: measurement errors and stellar jitter are
comparable to the signal, so for many stars the correct orbit cannot be
distinguished from an alias or pinned down within the 10\% amplitude
tolerance. We therefore stopped working on the search and focused on how
the method chooses among candidates and on the bias of the amplitude
estimate at low signal-to-noise. Full logs:
\url{https://github.com/zacharydoroshenko/Galileo-HypeCycle}
(\texttt{report/EXPERIMENT\_LOG.md}, \texttt{results/}).

\begin{table}[ht]
\caption{Decision trail. Run IDs match \texttt{report/EXPERIMENT\_LOG.md}.
``Proxy'' sets are built from TRAIN stars only (no simulated stars).}
\label{tab:trajectory}
\centering\small
\begin{tabularx}{\textwidth}{@{}lXXXX@{}}
\toprule
Run ID & Why & Change & Result & Next decision \\
\midrule
R01 & Need an exact copy of the baseline to measure changes against; our own solver was $>$5$\times$ slower & Reimplement the official baseline as \texttt{rv\_v2} with all switches off &
  Matches published validation to 4 decimals (0.4715 / 0.5975 / 0.316 / 0.375) & Build every change as a switch on \texttt{rv\_v2} \\
R02--03 & Task card: a perfect search adds only +0.0075, so test candidate choice and noise model & Ablation of 5 switches on 300-star train proxies; paired confirmation of \texttt{occam\_weight=0.5} &
  Occam best (mean 0.493$\to$0.503, no extra cost); confirmed on held-out train stars: std +17/$-$5 ($p$=0.017), sparse15 +11/$-$0 ($p$=0.001) & Keep Occam 0.5; drop yearly, peaks5, bounds \\
R04--05 & Sanity check before spending an upload & One-time validation check of Occam only; then a simulator-faithful sparse proxy (old proxy spans $\sim$4.5\% too long) &
  Validation gain far smaller than on proxies (in-dist +0.35, sparse +0.1, high-e $-$0.7 pts); faithful proxy reproduces the real sparse baseline exactly (0.375; old proxy 0.332) & Use the faithful proxy for all sparse decisions \\
R06--08, R13 & Test four specific hypotheses (Occam hurts eccentric fits; search misses optima; jitter bias; per-star noise for SNR) & \texttt{occam\_mode="pk"}; extra starts, $e$=0.92 grid, wider window; jitter fit with $N/(N-6)$; residual-noise SNR &
  None helped: pk std +2/$-$11 ($p$=0.02); search changes moved 0 of 1,000 std stars; jitter std +14/$-$28; residual SNR not significant ($p\ge$0.17) & All refuted; search is not the bottleneck \\
R09--11 & 26\% of baseline failures have the right $P$ but $K$ off by $>$10\%; ML $K$ is biased high at low SNR & $K/\sqrt{1+c/\mathrm{SNR}^2}$, $c$ fitted on 6,000 disjoint train stars; refit on augmented train variants &
  vs.\ Occam only: +70/$-$39 on 4,000 held-out std stars ($p\approx$0.004); same $c$ (6.93) on std and 15-obs stars & Final model = Occam 0.5 + $c$=6.95 \\
R12 & Confirm the final model before uploading & One-time validation check + official \texttt{headroom.py} &
  0.47475 / 0.6035 / 0.311 / 0.381 (88.6\% of the Fisher benchmark; baseline 88.0\%) & Upload as submission 16 (hidden test 0.4725); freeze \\
\bottomrule
\end{tabularx}
\end{table}

% ===========================================================================
\section{Final model \& comparisons}
% Owner: ZACH (method paragraph), PRISHA (tables, figure, binding info)
% NOTES FOR ZACH (method facts; write the paragraph yourself):
%  - candidates: GLS periodogram (1/1500 to 1/0.9 c/d, step 1e-4), top-3 peaks + 1-day aliases |1-f| (baseline)
%  - per candidate: grid of 396 (e, omega, M0) points, K and gamma solved by weighted linear LS; 3 best seed
%    scipy least_squares on (lnP, lnK, sqrt(e)cos w, sqrt(e)sin w, M0, gamma), freq kept within +-1.5 peak widths
%  - OURS: choose among fits by chi^2/2 + 0.5 * 0.5 * log det(J^T J) (Laplace/Occam factor, weight 0.5 chosen on train)
%  - OURS (if kept): K_out = K / sqrt(1 + c/SNR^2), SNR = K sqrt(N)/sig_eff, c fitted on train stars
%  - kind=algorithm; train_command only bakes the settings (+P/K bounds read from train labels); 0.31 s/star, 161 s for 4,000 stars at 14 workers
% DRAFTED FROM THE LOG BY THE CODING ASSISTANT: Zach, reread and reword in your own words.
Our final model keeps the official baseline's candidate search and changes only how a candidate is chosen and how
$K$ is reported. Candidates are the three highest peaks of a generalised Lomb--Scargle periodogram~\cite{zechmeister2009} plus their 1-day
aliases; for each, 396 $(e,\omega,M_0)$ grid points with $K$ and $\gamma$ solved by weighted linear least squares seed
a bounded six-parameter least-squares fit. Instead of the lowest $\chi^2$, we choose the fit minimising
$\chi^2/2 + w\cdot\tfrac12\log\det(J^\top J)$, a Laplace (Occam) factor~\cite{mackay2003} that prefers well-determined solutions
when the data are ambiguous, with $w=0.5$ selected on held-out train stars. We then correct the low-SNR upward bias~\cite{shen2008} of
the maximum-likelihood amplitude, $K_\mathrm{out}=K/\sqrt{1+c/\mathrm{SNR}^2}$ with $\mathrm{SNR}=K\sqrt{N}/\sigma_\mathrm{eff}$,
where $c=6.95$ was fitted on 6,000 train stars disjoint from all evaluation sets. The submission is \texttt:
\texttt{train\_command} writes the plugin with these two constants (and the train-label $P$/$K$ ranges, unused) and
learns no weights; it runs at 0.31\,CPU-s per star (161\,s for 4,000 stars with 14 workers).

\textbf{Binding.} Final submission ID: 16. ZIP SHA-256:
\texttt{1a4b4e94\ldots e2662} (full value in \texttt{report/SUBMISSIONS.md}). Git commit: 5c1a5b9. Splits: selection on TRAIN-derived proxies only; validation used once
as a held-out check; the leaderboard score is on the organizers' hidden test
set. Health gates (finite outputs, $0\le e<1$, positive $P,K$, dispersion of
$\log_{10}P \ge 0.05$ dex): all passed (local smoke tests on train and validation; the portal scored submission 16, which requires the blocking gates to pass).

\begin{table}[ht]
\caption{Comparable measurements (\texttt{success\_rate}, higher is better).
Validation numbers are a one-time held-out check, not a selection target.
Not-measured values are marked n/m.}
\label{tab:comparison}
\centering\small
\begin{tabularx}{\textwidth}{@{}lXccccc@{}}
\toprule
Method & Split & Overall & In-dist & High-e & Sparse & CPU-s/star \\
\midrule
Official baseline & validation & 0.4715 & 0.5975 & 0.316 & 0.375 & 0.24 \\
Official baseline & hidden test & n/a & n/a & n/a & n/a & -- \\
Final model & validation & 0.4748 & 0.6035 & 0.311 & 0.381 & 0.31 \\
Final model & hidden test & 0.4725 & n/a & n/a & n/a & -- \\
Alternative: Occam only & validation & 0.4718 & 0.601 & 0.309 & 0.376 & 0.31 \\
Alternative: $K$-debias only & validation & 0.4735 & 0.599 & 0.318 & 0.378 & 0.31 \\
Fisher ceiling & validation & 0.536 & 0.657 & 0.357 & 0.472 & -- \\
\bottomrule
\end{tabularx}
\end{table}

\begin{figure}[ht]
\centering
\begin{minipage}{0.49\linewidth}
\figorbox{figs/success_vs_ceiling.pdf}{Success by group: baseline, final, Fisher ceiling}
\end{minipage}\hfill
\begin{minipage}{0.49\linewidth}
\figorbox{figs/failure_breakdown.pdf}{Failure classes (headroom.py): K-off, 1-day alias, yearly alias, harmonic, other}
\end{minipage}
\caption{\fillin{PRISHA}{caption: what the reader should take away}}
\label{fig:main}
\end{figure}

% ===========================================================================
\section{Key insights}
% Owner: BOTH. Max three. Claim / evidence / limitation. Our own words.
\textbf{Insight 1.} \emph{Claim:} \fillin{BOTH}{}
\emph{Evidence:} \fillin{BOTH}{run/table/figure}
\emph{Limitation:} \fillin{BOTH}{}

\textbf{Insight 2.} \emph{Claim:} \fillin{BOTH}{candidate topic: gains on
train proxies (R03) shrank on validation (R04); why proxies mislead}
\emph{Evidence:} \fillin{BOTH}{}
\emph{Limitation:} \fillin{BOTH}{}

\textbf{Insight 3 (optional).} \emph{Claim:} \fillin{BOTH}{}
\emph{Evidence:} \fillin{BOTH}{}
\emph{Limitation:} \fillin{BOTH}{}

% ===========================================================================
\section{Limitations \& next experiment}
% Owner: PRISHA drafts, ZACH reviews
\fillin{PRISHA}{Most consequential uncertainty or failure mode (e.g. no train
star has $e>0.6$, so the high-e group cannot be selected on; simulated white
Gaussian noise; planet-free stars not scored). Then one specific next
experiment, the competing explanations it would separate, and how each
outcome would change our conclusion.}

% ===========================================================================
\section{Reproducibility, contributions \& AI use}
% Owner: PRISHA
\textbf{Reproduce.} \texttt{train\_command}:
\texttt{python train/train\_v2.py --data\_dir "\$DATA\_DIR" --seed "\$SEED" --out final\_model.py}.
Environment: SIDERIUS v1.1.0, Python 3.12, NumPy, SciPy (no other
dependencies). \texttt, seed 0 (the algorithm draws no
random numbers; the seed is recorded for the interface);
\texttt{initial\_algorithm.py} is the unmodified switchable source, and the
final plugin bakes in the selected settings. Data use: only
\texttt{train\_inputs.npy}/\texttt{train\_targets.npy} for selection; no
simulated stars; validation used once as a check. Code and complete logs:
\url{https://github.com/zacharydoroshenko/Galileo-HypeCycle} at commit 5c1a5b9 (submitted code; later commits add analysis only).
\fillin{PRISHA}{list any evidence that is unavailable}

\textbf{Contributions.} Zachary Doroshenko: \fillin{ZACH}{}.
Prisha Maiti: \fillin{PRISHA}{}.

\textbf{AI use.} \fillin{BOTH}{Which tools (Claude Code as coding agent; we did
/ did not use SIDERIUS capability calls), what they did (code, data analysis
scripts, organising logs, proofreading), and how outputs were verified
(metric reproduces the official baseline to 4 decimals; paired tests on
held-out train stars; local packaging checks). State that the hypotheses,
insights and limitations are our own.}

% OPTIONAL: delete this entire section if we do not find a reproducible loophole.
% It counts toward the four-page body.
\section*{Optional: Reward hacking discovery}
% FACTS (results/loophole_demo.json, tools/loophole_demo.py; LOCAL ONLY, never uploaded). Explanation, fix and
% false-positive discussion must be written by us (README: "Your team writes ... loophole explanations").
%  - e is not in the success criterion: replacing every e by 0.1 leaves success unchanged and passes all health gates,
%    while median |de| rises 0.059 -> 0.152 on validation; in the high_ecc group 0.065 -> 0.648 (10x), in-dist 0.042 -> 0.089.
%  - The K box |K'-K|/K < 0.1 is asymmetric in log space ([-0.0458, +0.0414] dex); K x 0.995 recentres it:
%    +14/-10 on validation, +19/-13 on 6,000 train stars (n.s., +0.1 pt). It cannot explain our K-debias gain
%    (+70/-39, and the threshold-free median rel K error also drops 0.0790 -> 0.0751).
%  - Possible fixes: success criterion in log space |log K'/K| < log 1.1; an e-accuracy gate or e in the criterion.
%    False positive to discuss: at low SNR e is genuinely unconstrained, so an e gate would penalise honest methods.
\begin{table}[ht]
\caption{Metric audit on validation (local; not a submitted model). The ``loophole'' variant changes no physics.}
\centering\small
\begin{tabular}{@{}lcccc@{}}
\toprule
Predictions & Success & vs.\ final & Median $|\Delta e|$ & Median $|\Delta K|/K$ \\
\midrule
Final model (submission 16) & 0.4748 & -- & 0.059 & 0.0751 \\
$e \to 0.1$ for every star & 0.4748 & +0/$-$0 & 0.152 & 0.0751 \\
$K \to 0.995\,K$ (recentre box) & 0.4758 & +14/$-$10 & 0.059 & 0.0742 \\
Both (``loophole model'') & 0.4758 & +14/$-$10 & 0.152 & 0.0742 \\
\bottomrule
\end{tabular}
\end{table}
\fillin{PRISHA}{Only if reproducible. Scientific goal vs metric incentive;
triggering condition; minimal reproducible example; evidence that a high
score coexists with scientific failure; proposed fix and its false positives.
Leads from the task card: planet-free stars are never scored (false
detections are free); public simulator can remove the out-of-distribution
premise; sparse stars are identifiable from their 15-point mask.}

% References may extend beyond the four-page body. Do not add narrative appendices.
\clearpage
\begin{thebibliography}{9}
\bibitem{galileo2026} Galileo Hackathon organizers. \texttt{exoplanet-orbits} task package (TASK\_CARD, simulator
\texttt{generate\_rv\_data.py} v2, official baseline, \texttt{analysis/headroom.py}) and data. Zenodo,
doi:10.5281/zenodo.24361910, 2026.
\bibitem{challenge} Galileo Hackathon 2026. Recover exoplanet orbits from
sparse radial-velocity data (Track~2 challenge \texttt{exoplanet-orbits}).
\url{https://siderius.scitra.org/hackathon/challenges/exoplanet-orbits}.
\bibitem{zechmeister2009} M.~Zechmeister and M.~K\"urster. The generalised
Lomb--Scargle periodogram. \emph{A\&A} 496:577--584, 2009.
\bibitem{dawson2010} R.~I.~Dawson and D.~C.~Fabrycky. Radial velocity planets de-aliased. \emph{ApJ} 722:937, 2010.
\bibitem{mackay2003} D.~J.~C.~MacKay. \emph{Information Theory, Inference, and Learning Algorithms}, ch.~28
(Model comparison and Occam's razor). Cambridge University Press, 2003.
\bibitem{shen2008} Y.~Shen and E.~L.~Turner. On the eccentricity distribution of exoplanets from radial velocity
surveys. \emph{ApJ} 685:553, 2008.
\bibitem{pricewhelan2017} A.~M.~Price-Whelan et al. The Joker: a custom Monte Carlo sampler for binary-star and
exoplanet radial velocity data. \emph{ApJ} 837:20, 2017.
\bibitem{kipping2013} D.~M.~Kipping. Parametrizing the exoplanet eccentricity distribution with the beta
distribution. \emph{MNRAS} 434:L51, 2013.
\end{thebibliography}
\end{document}
